\documentclass[11pt]{article}
\usepackage{fullpage}
\usepackage{amsmath, multirow}
\usepackage{amssymb}
\usepackage{amsthm}
\usepackage{xcolor}

\newcommand{\dist}{\mathrm{\bf dist}}
\pagestyle{empty}

%\parskip .125in

\begin{document}
\begin{center}
{\bf Intro to Computational Math, Fall 2025\\
Homework 5 Sample Solutions}
\end{center}

{\bf Question 1.}

\paragraph{3.1.3}
Let the weekly earnings be $y_i>0$ at times $t_i=1,\dots,18$. Model $y(t)=a e^{b t}$. Taking logs gives
\[
\log y_i = c + b t_i,\quad c=\log a.
\]
With design matrix \(X=[\mathbf{1}\ t]\) and vector \(Y=(\log y_i)\), the least-squares solution is
\[
\begin{bmatrix}c\\ b\end{bmatrix}=(X^{T}X)^{-1}X^{T}Y.
\]
Using the given data yields
\[
b\approx -0.34759,\qquad a=\exp(c)\approx 1.273\times 10^{8}.
\]
Thus the fit is \(y(t)\approx 1.273\times 10^{8}\, e^{-0.34759 t}\). The multiplicative week-to-week factor is \(e^{b}\approx 0.7064\).

\medskip
\noindent\textbf{Python:}
\begin{verbatim}
import numpy as np
import matplotlib.pyplot as plt

y = np.array([189923838, 79406327, 46230374, 26830921, 18804290,
              13822428, 7474688, 6129424, 4377675, 3764963, 2426574,
              1713298, 1426102, 1031985, 694947, 518242, 460578, 317909],
             dtype=float)

t = np.arange(1, len(y)+1)            # weeks 1..18
X = np.vstack([np.ones_like(t), t]).T # [1, t]
coef = np.linalg.lstsq(X, np.log(y), rcond=None)[0]
c, b = coef[0], coef[1]
a = np.exp(c)

print("a =", a, "   b =", b, "   exp(b) =", np.exp(b))

# Linear axes plot
tt = np.linspace(1, len(y), 300)
yfit = a * np.exp(b * tt)
plt.figure()
plt.plot(t, y, "o", label="data")
plt.plot(tt, yfit, "-", label="fit y=a*exp(b t)")
plt.xlabel("week")
plt.ylabel("weekly gross (USD)")
plt.legend()
plt.title("Hunger Games weekly box office: linear axes")
plt.show()

# Semilogy plot (log scale on y)
plt.figure()
plt.semilogy(t, y, "o", label="data")
plt.semilogy(tt, yfit, "-", label="fit y=a*exp(b t)")
plt.xlabel("week")
plt.ylabel("weekly gross (USD)")
plt.legend()
plt.title("Hunger Games weekly box office: semilogy")
plt.show()
\end{verbatim}

\paragraph{3.1.7}
Model the orbital period \(\tau\) versus mean distance \(R\) by \(\tau=c R^{\alpha}\) with \(\tau>0\), \(R>0\). Taking logs gives
\[
\log \tau_i = \log c + \alpha \log R_i .
\]
With \(x_i=\log R_i\), \(Y_i=\log \tau_i\), and \(X=[\mathbf{1}\ x]\), the least-squares solution
\[
\begin{bmatrix}\log c\\ \alpha\end{bmatrix}=(X^{T}X)^{-1}X^{T}Y
\]
from the table data is
\[
\alpha \approx 1.49783,\qquad c\approx 0.20178 .
\]
Hence \(\tau \approx 0.20178\, R^{1.49783}\). The exponent is extremely close to the simple rational value \(\alpha=3/2\), consistent with Kepler's third law.

\medskip
\noindent\textbf{Python:}
\begin{verbatim}
import numpy as np
import matplotlib.pyplot as plt

R = np.array([57.9, 108.2, 149.6, 228.0, 778.5, 1432.0, 2867.0, 4515.0],
             dtype=float)   # Mkm
tau = np.array([88.0, 224.7, 365.2, 687.0, 4331.0, 10750.0, 30590.0, 59800.0],
               dtype=float) # days

X = np.vstack([np.ones_like(R), np.log(R)]).T
coef = np.linalg.lstsq(X, np.log(tau), rcond=None)[0]
lnc, alpha = coef[0], coef[1]
c = np.exp(lnc)
print("alpha =", alpha, "   c =", c)

# Log-log plot with fitted line
xx = np.linspace(np.min(np.log(R)), np.max(np.log(R)), 200)
yy = lnc + alpha * xx
plt.figure()
plt.plot(np.log(R), np.log(tau), "o", label="data")
plt.plot(xx, yy, "-", label="fit log tau = log c + alpha log R")
plt.xlabel("log R (Mkm)")
plt.ylabel("log tau (days)")
plt.legend()
plt.title("Planetary periods vs distance: log-log least squares")
plt.show()
\end{verbatim}


\newpage



{\bf Question 2.}

\paragraph{3.2.1}
The normal equations are \(A^{T}A x = A^{T}b\). With
\[
A=\begin{bmatrix}2&-1\\0&1\\-2&2\end{bmatrix},\quad
b=\begin{bmatrix}1\\-5\\6\end{bmatrix},
\]
compute
\[
A^{T}A=
\begin{bmatrix}
8&-6\\
-6&6
\end{bmatrix},
\qquad
A^{T}b=
\begin{bmatrix}
-10\\
6
\end{bmatrix}.
\]
Solve
\[
\begin{bmatrix}
8&-6\\
-6&6
\end{bmatrix}
\begin{bmatrix}x_1\\x_2\end{bmatrix}
=
\begin{bmatrix}-10\\6\end{bmatrix}.
\]
Adding the two equations gives \(2x_1=-4\) so \(x_1=-2\). Substituting into \(-6x_1+6x_2=6\) yields \(12+6x_2=6\), hence \(x_2=-1\). Therefore
\[
x^{\ast}=\begin{bmatrix}-2\\-1\end{bmatrix}.
\]


\newpage



{\bf Question 3.}

\paragraph{(i)}
For \(z=\begin{bmatrix}6\\ -2\\ 9\end{bmatrix}\) we have \(\|z\|=\sqrt{36+4+81}=11\). Then
\[
v=\|z\|e_1-z=\begin{bmatrix}11\\0\\0\end{bmatrix}-\begin{bmatrix}6\\-2\\9\end{bmatrix}
=\begin{bmatrix}5\\2\\-9\end{bmatrix},
\qquad
v^{T}v=25+4+81=110.
\]
Therefore
\[
P=I-\frac{2}{v^{T}v}vv^{T}
=I-\frac{1}{55}\begin{bmatrix}
25&10&-45\\
10&4&-18\\
-45&-18&81
\end{bmatrix}
=
\begin{bmatrix}
6/11&-2/11&9/11\\
-2/11&51/55&18/55\\
9/11&18/55&-26/55
\end{bmatrix}.
\]
A direct multiplication gives
\[
Pz=\begin{bmatrix}11\\0\\0\end{bmatrix}.
\]

\paragraph{(ii)}
Let \(v\neq 0\) and define \(P=I-2\,vv^{T}/(v^{T}v)\). Since \(vv^{T}\) is symmetric, \(P^{T}=I-2\,vv^{T}/(v^{T}v)=P\), so \(P\) is symmetric.
Let \(u=v/\|v\|\). Then \(P=I-2uu^{T}\) and \(u^{T}u=1\). Compute
\[
P^{2}=(I-2uu^{T})^{2}=I-4uu^{T}+4u(u^{T}u)u^{T}=I-4uu^{T}+4uu^{T}=I.
\]
Hence \(P^{T}P=P^{2}=I\) and \(P\) is orthogonal.

\paragraph{(iii)}
Take \(z=\begin{bmatrix}10^{8}\\1\\1\end{bmatrix}\). Then \(\|z\|=\sqrt{10^{16}+2}\) is extremely close to \(10^{8}\), so
\[
v_1=\|z\|-z_1=\sqrt{10^{16}+2}-10^{8}
\]
is the subtraction of two nearly equal positive numbers, causing loss of significant digits in finite precision. Thus \(v=\|z\|e_1-z\) suffers subtractive cancellation.

\paragraph{(iv)}
Let
\[
v_{+}=\|z\|e_1-z,\qquad v_{-}=-\|z\|e_1-z.
\]
A short calculation gives
\[
\|v_{+}\|^{2}=2\|z\|(\|z\|-z_1)=2\|z\|^{2}\Bigl(1-\frac{z_1}{\|z\|}\Bigr),
\qquad
\|v_{-}\|^{2}=2\|z\|(\|z\|+z_1)=2\|z\|^{2}\Bigl(1+\frac{z_1}{\|z\|}\Bigr).
\]
Hence
\[
\Bigl(\frac{\|v_{+}\|}{\|z\|}\Bigr)^{2}=2\Bigl(1-\frac{z_1}{\|z\|}\Bigr),\qquad
\Bigl(\frac{\|v_{-}\|}{\|z\|}\Bigr)^{2}=2\Bigl(1+\frac{z_1}{\|z\|}\Bigr).
\]
Since \(z_1/\|z\|\in[-1,1]\), at least one of \(1\pm z_1/\|z\|\) is at least \(1\). Therefore at least one of the two ratios is at least \(\sqrt{2}\), in particular at least \(1\). A practical rule is to choose the sign matching \(\operatorname{sign}(z_1)\), which yields the larger of the two norms and avoids subtractive cancellation in the first component.




\end{document}
